\section{Discussion and Qualification Limits}
\label{sec:discussion}
\subsection{What the Evidence Establishes}
The numerical tests support the specified discrete transformation, including cases where an intuitive floating-point replacement would differ by signed rounding or limiting order. Renderer pixel comparisons establish consistency with the preserved optical model. Source/pose tests show that accepted pass state is frozen and that invalid associations fall back to static correction. Routing establishes that the constrained implementation meets the reported timing checks on the selected device. These are complementary results, with distinct oracles and scopes.

The saved software outputs provide static alignment comparisons and paired physical-marker orientation measurements. The latter show lower image-angle variation for the saved inertial output in two selected intervals, without an independent gravity reference. These experiments do not test the current FPGA pixel path, which intentionally retains older optical maps, or the current live inertial loop. Likewise, historical FPS and latency measurements demonstrate behavior of their recorded builds. Their utility for this design is to identify repeated-camera frames, buffer publication delay, and bandwidth pressure that a purely geometric validation would miss.

\subsection{Limits of the Motion and Sensor Model}
Only rig roll and pitch enter the live rotation. Yaw heading stabilization, depth-dependent translation, rolling-shutter compensation, and inter-camera radiometric correction are outside this model. The affine source-map correction is an approximation whose acceptable error depends on field of view, tilt, crop, and scene depth. Saturation retains memory safety but can leave residual geometric error. A one-pass global limiter is not a proof that every camera output lies within the renderer envelope when mounting trims are nonzero.

The selected mounting profile and the preserved optical maps have separate provenance. Changing optical focal parameters, pitch scale, or crop requires regenerating and verifying the maps, signed shift widths, and capture envelope together. Dynamic UI roll/pitch trims are supported; saving a yaw value does not change the current RTL's fixed mounting yaw. A stored calibration marker prevents first-start defaults from overwriting subsequent UI edits, but a programmer mass erase removes that state.

\subsection{Dynamic Validation Still Required}
The authors' qualitative report of live functionality is consistent with the implementation intent, but does not supply an error distribution or a reproducible image/firmware identity. A quantitative test should record the FPGA bitstream/probe hashes, MCU image, actual trim banks, camera format, exposure, trigger delay, DNR state, and measured source and output rates. A fixed rig motion profile and an independent orientation reference should be captured together with source-exposure and displayed-frame identities. Report both leveling and temporal smoothness because suppressing low-frequency gravity error and smoothing relative image motion are different objectives.

At minimum, compare stabilization Off/On on identical motions; test receipt-time versus frame-associated pose; exercise the static fallback and calibration update boundaries; and repeat under DDR stalls and each sensor family. Measure angle error, valid area, seam disagreement, distinct-frame throughput, and event-to-event latency with distributions across repeated runs. Controlled exposure timing is necessary when comparing camera synchronization: raster-start alignment alone can remain good while exposure or ISP latency changes. A current-build sensor-to-display latency claim requires the complete acquisition/processing/scanout path to be observed, rather than adding nominal half-frame terms to a diagnostic timestamp.

\subsection{Reproducibility}
The revision archive records source revisions, calibration and optical-asset hashes, derived comparison tables, build reports, and the scripts used to produce the evidence summaries. The video audit retains all 2700 method/frame rows across the two marker intervals, including missing detections, marker IDs, corners, frame hashes, and common-support flags. Separate audits reproduce the synthetic-horizon and IMU diagnostics while recording their measurement limitations. Local source datasets and captures are identified by path and hash; this is provenance, not a claim that every original capture is publicly released. The main tables state excluded failures, tested subsets, and historical build identities. No learning-based IR comparison or independent current-hardware dynamic accuracy measurement is implied.
